\chapter{Ideale und Filter im Verband}
\label{b45guide:ch:ideals}

\section{Operationen bestimmen die Ideale}

\begin{definition}[Ideal und Filter]
\label{b45guide:def:ideal}
Ein \emph{Ideal} eines supremal gelesenen Halbverbandes $S$
ist eine Teilmenge $I\subseteq S$, für die
\[
 x\in I,\ y\leq x\Rightarrow y\in I,
 \qquad x,y\in I\Rightarrow x\vee y\in I
\]
gilt. Ein \emph{Filter} eines infimal gelesenen Halbverbandes
ist entsprechend nach oben abgeschlossen und unter $\wedge$
abgeschlossen.
Hier werden das leere Ideal und der leere Filter zugelassen.
Ein \emph{$0$-Ideal} enthält zusätzlich die benannte Null;
ein \emph{$1$-Filter} enthält die benannte Eins.
\end{definition}

\begin{theorem}[Algebraisches Idealkriterium]
\label{b45guide:thm:ideal-test}
Für eine Teilmenge $I$ eines Halbverbandes $(S,\vee)$ sind äquivalent:
\begin{enumerate}
\item $I$ ist ein Ideal.
\item Für alle $x,y\in S$ gilt
\[
 x\vee y\in I\quad\Longleftrightarrow\quad x\in I\ \land\ y\in I.
\]
\item Die Abbildung $\chi_I:S\to\mathbf2$, die auf $I$ den Wert $0$
und außerhalb von $I$ den Wert $1$ annimmt, erhält $\vee$.
\end{enumerate}
Besitzt $S$ eine benannte Null, ist $I$ genau dann ein $0$-Ideal,
wenn $\chi_I$ zusätzlich diese Null erhält. \cite{smith}
\end{theorem}

\begin{definition}[Hauptideal und erzeugtes Ideal]
\label{b45guide:def:principal}
Für $a\in S$ ist
\[
 \mathord\downarrow a=\{x\in S\mid x\leq a\}
\]
das \emph{Hauptideal} von $a$.
Das kleinste Ideal, das $X\subseteq S$ enthält, heißt
$\operatorname{Id}(X)$. Dual bezeichnet
$\mathord\uparrow a=\{x\in S\mid a\leq x\}$ den Hauptfilter,
und $\operatorname{Fil}(X)$ den von $X$ erzeugten Filter.
\end{definition}

\begin{theorem}[Formel für erzeugte Ideale]
\label{b45guide:thm:generated-ideal}
Im Halbverband $(S,\vee)$ gilt
\[
 \operatorname{Id}(X)=
 \{z\in S\mid
     z\leq x_1\vee\cdots\vee x_n
     \text{ für ein }n\geq1\text{ und }x_1,\ldots,x_n\in X\}.
\]
Insbesondere ist $\operatorname{Id}(\varnothing)=\varnothing$.
Für ein $0$-Ideal wird auch der leere Term $\bigvee\varnothing=0$
zugelassen; das von der leeren Menge erzeugte $0$-Ideal ist $\{0\}$.
Duale Formeln gelten für Filter mit endlichen
$\wedge$-Verknüpfungen. \cite{nation-ideals}
\end{theorem}

\begin{theorem}[Endliche Erzeugung liefert Hauptideale]
\label{b45guide:thm:finite-ideal}
Für jede nichtleere endliche Menge $X$ in einem Halbverband gilt
\[
 \operatorname{Id}(X)=\mathord\downarrow\!\bigl(\bigvee X\bigr).
\]
Daher sind die nichtleer endlich erzeugten Ideale genau die
Hauptideale. Jedes nichtleere Ideal eines endlichen Halbverbandes
ist ein Hauptideal. Dual sind die nichtleer endlich erzeugten
Filter genau die Hauptfilter.
\end{theorem}

\section{Der Verband der Ideale}

\begin{theorem}[Idealverband und kanonische Darstellung]
\label{b45guide:thm:ideal-lattice}
Sei $L$ ein Verband, und $\operatorname{Idl}(L)$ die Menge seiner
Ideale einschließlich des leeren Ideals. Die Operationen
\[
 I\wedge_I J=I\cap J,\qquad
 I\vee_I J=\operatorname{Id}(I\cup J)
\]
machen $\operatorname{Idl}(L)$ zu einem Verband.
Für nichtleere $I,J$ gilt die konkrete Formel
\[
 I\vee_I J=\{x\in L\mid
               x\leq i\vee j\text{ für gewisse }i\in I,\ j\in J\}.
\]
Das leere Ideal ist neutrales Element für $\vee_I$.

Die Hauptidealabbildung $a\mapsto\mathord\downarrow a$ ist eine
Verbandseinbettung:
\[
\begin{aligned}
 \mathord\downarrow(a\wedge b)
   &=(\mathord\downarrow a)\cap(\mathord\downarrow b),\\
 \mathord\downarrow(a\vee b)
   &=(\mathord\downarrow a)\vee_I(\mathord\downarrow b).
\end{aligned}
\]
Dabei ist $\vee_I$ die Idealbildung der Vereinigung, im Allgemeinen
nicht die bloße Mengenvereinigung. \cite{nation-ideals}
\end{theorem}

\begin{theorem}[Hauptideale bei endlichen Verbänden]
Ist $L$ ein nichtleerer endlicher Verband, so ist die
Hauptidealabbildung ein Isomorphismus von $L$ auf den Verband seiner
\emph{nichtleeren} Ideale mit den eben definierten Operationen.
Der Verband $\operatorname{Idl}(L)$ mit zusätzlich leerem Ideal
besitzt genau $|L|+1$ Elemente.
\end{theorem}
